% !TEX root = ../my-thesis.tex
%
\chapter{Einleitung}
\label{sec:intro}

\section{Motivation}
\label{sec:intro:motivation}

\section{Problemstellung}
\label{sec:intro:problem}

\section{Ziel der Arbeit}
\label{sec:intro:goal}

\section{Aufbau der Arbeit}
\label{sec:intro:structure}

Die Arbeit ist wie folgt aufgebaut:


\textbf{Kapitel \ref{sec:basics}} \\[0.2em]
Dieses Kapitel stellt die begrifflichen Grundlagen vor, die für das Verständnis der weiteren Arbeit notwendig sind.

\textbf{Kapitel \ref{sec:related}} \\[0.2em]
Dieses Kapitel behandelt verwandte Arbeiten.

\textbf{Kapitel \ref{sec:classification}} \\[0.2em]
Hier wird ein Kriterienrahmen zur Klassifikation von Workflow-Management-Systemen vorgestellt.

\textbf{Kapitel \ref{sec:methodology}} \\[0.2em]
Dieses Kapitel beschreibt die methodische Vorgehensweise der Arbeit und die Auswahl der betrachteten Systeme.

\textbf{Kapitel \ref{sec:description}} \\[0.2em]
Hier erfolgt die bEschreibung und der Vergleich ausgewählter Workflow-Management-Systeme anhand der zuvor definierten Kriterien.

\textbf{Kapitel \ref{sec:practical}} \\[0.2em]
Dieses Kapitel behandelt die praktische Modellierung eines Beispiel-Workflows in ausgewählten Workflow-Management-Systemen.

\textbf{Kapitel \ref{sec:discussion}} \\[0.2em]
In diesem Kapitel werden die Ergebnisse des Vergleichs und der praktischen Untersuchung diskutiert.

\textbf{Kapitel \ref{sec:conclusion}} \\[0.2em]
Abschließend werden die wichtigsten Ergebnisse zusammengefasst.